\subsection{Concordance components from dimension six}

The stable presentation already computes the unstable group for
smoothable manifolds from dimension six. The reason is geometric:
positive suspension preserves both the first obstruction and the
complete framed second obstruction. The corrected geometric exact
sequence then identifies the whole group, including its extension
law. This argument uses Igusa's correction to the
Hatcher--Wagoner sequence; it does not discard the framing term or
choose a splitting.

\begin{theorem}\label{thm:smoothable-concordance-six}
Let $M$ be a compact connected smoothable PL manifold of dimension
$n\geq6$, possibly with boundary, and let $K$ be a finite simplicial
complex homotopy equivalent to $M$. With the bottom and the entire
side boundary fixed pointwise and the upper face free,
\[
                     \pi_0P_{\mathrm{PL}}(M)\cong E(K).
\]
PL components mean jointly PL isotopies. The same group presents
the stipulated definable concordance components of any chosen
finite model over any real closed field and in any fixed o-minimal
expansion. Finite disjoint unions give finite products, provided
every component is smoothable and has dimension at least six.

For a compact smooth manifold of dimension at least six, positive
suspension induces a group isomorphism
\[
 \sigma^+_*:\pi_0P_{\mathrm{Diff}}(M)
       \xrightarrow{\ \cong\ }\pi_0P_{\mathrm{Diff}}(M\times I).
\]
No smoothing or isotopy is prescribed on the boundary in advance;
the concordances themselves fix the side boundary. No assertion
about higher homotopy groups is made here.
\end{theorem}

\begin{proof}
For a connected space $Z$, set
\[
 A_Z=\operatorname{Wh}_1^+
       (\pi_1Z;\mathbb Z/2\oplus\pi_2Z),\qquad
 I_Z=\operatorname{im}(\chi_Z:\operatorname{Wh}_3(\pi_1Z)\to A_Z).
\]
Here $\chi_Z$ is the full corrected obstruction homomorphism.  Its
image is unchanged if its domain is replaced by the surjecting group
$K_3(\mathbb Z[\pi_1Z])$.
For a compact smooth $n$-manifold, $n\geq6$, the corrected geometric
sequence gives
\begin{equation}\label{eq:full-geometric-concordance-sequence}
 0\longrightarrow A_M/I_M\xrightarrow{\overline\theta_M}
 \pi_0P_{\mathrm{Diff}}(M)\xrightarrow{\Sigma_M}
 \operatorname{Wh}_2(\pi_1M)\longrightarrow0.
\end{equation}
This is Igusa's corrected \emph{unstable} geometric sequence
\cite[Section~8.a and Theorem~8.a.2, pp.~149--150]{Igusa1984}.
Chapter~8 assumes $\dim M\geq6$. The map $\chi_M$ includes both
the Postnikov contribution and the Grassmann, or framing,
contribution. In particular, we take the cokernel of the full map
\[
 K_3(\mathbb Z[\pi_1M])\longrightarrow
 \operatorname{Wh}_1^+(\pi_1M;\pi_2M)
 \oplus\operatorname{Wh}_1^+(\pi_1M;\mathbb Z/2).
\]
The fixed-boundary convention is specified on printed p.~104.
The sequence is also stated in
\cite[Theorem~3.1, p.~6]{Igusa2021Dimension3}; its functorial stable
form is \cite[Corollary~2.7, p.~250]{Jahren2010}. We use the
geometric sequence, rather than infer an unstable dimension range
from the stable formula.

Projection $M\times I\to M$ identifies the fundamental groups and their
modules $\pi_2$.  Functoriality of the full $\chi$ identifies both $A$
and $I$ for these two spaces.  To apply these identifications to the
geometric sequence, we check its two obstruction maps under suspension.

Normalize concordances to be the identity near the bottom and side
boundary and products near the top.  Collar straightening induces the
same components and works for paths of concordances
\cite[Part~I, Chapter~I, Section~5, p.~55]{HatcherWagoner1973}.
In the positive-suspension construction one applies a normalized
concordance on each of a fixed collection of polar rays and uses the
identity outside them.  This is the formula in
\cite[Section~4.1, p.~8]{Igusa2021Dimension3}.  Since the ray coordinates
are fixed independently of the concordance, it gives
\[
                    \sigma^+(gh)=\sigma^+(g)\sigma^+(h).
\]
In particular it induces a homomorphism on components, with the required
boundary restrictions.

For the second obstruction choose marked lens models in indices $3,4$
in every dimension $n\geq6$.  Hatcher and Wagoner's two-index theorem
permits $2\leq i\leq n-2$ for the one-parameter representatives
\cite[Part~I, Chapter~V, Theorem~3.1(a), p.~184;
Part~II, Section~2, p.~244]{HatcherWagoner1973}.
The marked-lens calculation identifies its group with
\[
 \frac{(\mathbb Z/2\oplus\pi_2M)[\pi_1M]}
      {(\mathbb Z/2\oplus\pi_2M)[1]}.
\]
For this specialization keep $k=3$. Then
\[
                         n-k\geq k=3
\]
for every $n\geq6$, exactly the inequality used in the proof of
\cite[Lemma~1.10, p.~9]{Igusa2021SecondObstructionI}.
The theorem there is stated with $k=\lfloor n/2\rfloor$; the
fixed-index specialization follows from its proof, the preceding
two-index representatives, and the general fixed-index suspension
spaces of \cite[Section~4.2, pp.~8--9]{Igusa2021Dimension3}.
Changing markings passes to $A_M$, and forgetting the lens data gives
$\theta_M$.  Use the stable sign convention
$\lambda^s_0=(-1)^k\lambda_0$.
Positive suspension keeps $k$ fixed and preserves the complete framed
invariant, including its $\mathbb Z/2$ coordinate
\cite[Section~4.2, pp.~8--9]{Igusa2021Dimension3}.  Thus, under the
identification of the groups $A$,
\[
                  \sigma^+_*\theta_M=\theta_{M\times I}.
\]
The first obstruction is also unchanged: this is the suspension remark
in \cite[Part~I, Chapter~IV, Section~1, p.~130]{HatcherWagoner1973},
based on the trajectory formula (5.2) on p.~60.  Therefore
\[
                    \Sigma_{M\times I}\sigma^+_*=\Sigma_M.
\]

We have obtained a commutative diagram of
\eqref{eq:full-geometric-concordance-sequence} for $M$ and $M\times I$
whose outside vertical arrows are isomorphisms and whose middle arrow
is $\sigma^+_*$.  Exactness proves that middle arrow is an isomorphism.
For injectivity, a class killed by suspension has zero first obstruction,
so lies in $A_M/I_M$; injectivity of the lower left map then kills it.
For surjectivity, first match a target class's first obstruction and then
correct the residual element using $A_M/I_M$.  This uses no splitting
of the extension.

All these statements concern compact manifolds with side boundary
fixed.  Igusa explicitly uses this convention on p.~1 and fixes the
projection near the full boundary in the setup on p.~3 of
\cite{Igusa2021SecondObstructionI}.  The handle and lens modifications
occur in the interior.  Thus the suspension argument may be iterated
on $M\times I^j$, with corners rounded, and identifies the unstable
smooth group with the stable group.  The countable presentation theorem
and the classical smooth-to-PL component comparison therefore give
$E(K)$ whenever $M$ is smoothable and $n\geq6$.

For clarity, the smooth--PL component comparison preserves
\emph{composition}. Let $u$ and $v$ be the forgetful homomorphisms
from the smooth and PL concordance component groups to the TOP
concordance component group. Theorems~6.1 and~6.2 of
\cite[pp.~41--43]{BurgheleaLashof1974II} give the required
component bijections and commute with these forgetful maps.
Consequently $v^{-1}u$ is the comparison group isomorphism.
This does not require a group of piecewise differentiable maps.

The relative PL--definable component comparison applies with
$M\times\{0\}\cup\partial M\times I$ as fixed subpolyhedron.
It preserves composition and transfers the presentation to every
stipulated real closed field model. Fixing the bottom prevents
permutation of connected components, proving the finite-product
assertion.
\end{proof}

\begin{corollary}\label{cor:pl-six-seven-components}
Every compact PL manifold whose connected components have
dimension six or seven satisfies the component presentation in
Theorem~\ref{thm:smoothable-concordance-six}, without an additional
smoothability assumption. The same is true of its stipulated
definable models.
\end{corollary}

\begin{proof}
The smoothing classification uses the fibration
$BO\to BPL$ with fibre $PL/O$
\cite[Part~II, Theorem~4.1, p.~108, and Theorem~5.3,
pp.~122--123]{HirschMazur1974}. The fibre is $6$-connected
\cite[Part~II, p.~123]{HirschMazur1974}. An obstruction to extending
a lift across a $j$-cell lies in $\pi_{j-1}(PL/O)$.
All these groups vanish for $j\leq7$, so the stable PL tangent
classifying map of a manifold of dimension at most seven admits
a lift. The smoothing theorem turns this lift into a compatible
smoothing.

For a manifold with boundary, first apply the boundaryless
classification to its interior, which has the same homotopy type.
A PL collar gives a compact interior copy of the original manifold.
Its locally flat boundary may be moved to a smooth hypersurface by
\cite[Part~I, Theorem~7.4, pp.~37--38]{HirschMazur1974}.
The enclosed compact smooth domain is PL homeomorphic to the
original manifold and supplies a compatible smoothing, including
the boundary. There was no prescribed boundary smoothing to extend.
The preceding theorem now applies.
\end{proof}

\begin{lemma}\label{lem:three-skeleton-cycle-presentation}
For every connected finite simplicial complex $K$, inclusion of
the three-skeleton induces a group isomorphism
\[
                              E(K^{(3)})\cong E(K).
\]
\end{lemma}

\begin{proof}
The functorial corrected stable sequence is
\[
 0\longrightarrow A_Z/I_Z\longrightarrow
 \pi_2\operatorname{Wh}^{\mathrm{Diff}}(Z)\longrightarrow
 \operatorname{Wh}_2(\pi_1Z)\longrightarrow0,
\]
with $A_Z$ and $I_Z$ as in the proof above
\cite[Corollary~2.7, p.~250]{Jahren2010}.
The inclusion $K^{(3)}\to K$ induces isomorphisms on $\pi_1$
and on its module $\pi_2$, since the remaining cells have
dimension at least four. It therefore induces isomorphisms on
$\operatorname{Wh}_3$, $A$, and $\operatorname{Wh}_2$.
Naturality of the complete $\chi$ identifies its image subgroups
$I$; no assertion that these images are determined by abstract
groups without their Postnikov data is needed. The diagram of
short exact sequences then identifies the middle terms.
The cycle-presentation identification
$E(Z)\cong\pi_2\operatorname{Wh}^{\mathrm{Diff}}(Z)$ gives the
claimed isomorphism.
\end{proof}
